% !TeX root = ../main.tex
\chapter{System Evaluation Framework}
\label{chap:system-evaluation}

\section{Purpose and Evaluation Scope}

At the current stage of the project, the implementation has not yet reached
a complete production-ready state. Therefore, this chapter does not present
empirical performance numbers or final experimental results. Instead, it
defines an evaluation framework that can be used once the implementation is
complete enough for pilot testing. The goal is to make the later evaluation
measurable, reproducible, and aligned with the requirements defined in
Chapter~\ref{chap:khao-sat-phan-tich-he-thong}.

The evaluation framework answers five questions:

\begin{enumerate}
    \item Does the system satisfy the core functional requirements for
          instructors, learners, administrators, and Canvas LMS integration?
    \item Does the architecture preserve correct authorization, data
          isolation, and consistency across PostgreSQL, Qdrant, Neo4j, Redis,
          and MinIO?
    \item Does retrieval return relevant and properly scoped learning
          materials for search, quiz generation, and AI Tutor answers?
    \item Are generated lesson cards, flashcards, quiz items, and chatbot
          responses pedagogically useful, factually grounded, and aligned with
          Learning Outcomes?
    \item Can the system meet practical non-functional expectations for
          latency, throughput, reliability, observability, and cost control in
          a demo or pilot deployment?
\end{enumerate}

Because the project is currently design-oriented, all criteria in this chapter
are expressed as \textbf{planned evaluation criteria}. A later implementation
phase should replace the planned protocol with measured results, datasets,
tool versions, and experiment logs.

\section{Evaluation Principles}

The evaluation is designed around four principles:

\begin{itemize}
    \item \textbf{Traceability.} Each evaluation item must map back to a user
          story, functional requirement, or non-functional requirement from
          Chapter~\ref{chap:khao-sat-phan-tich-he-thong}.
    \item \textbf{Reproducibility.} Each experiment must record dataset
          version, prompt version, model version, retrieval configuration,
          database schema version, and deployment configuration.
    \item \textbf{Separation of concerns.} Functional correctness, retrieval
          quality, generated-content quality, security, and performance are
          evaluated separately before being considered together.
    \item \textbf{No fabricated results.} When the system is not yet fully
          implemented, the project only states evaluation methods and expected
          acceptance criteria, not unmeasured success numbers.
\end{itemize}

\section{Evaluation Dimensions}

The system should be evaluated across seven dimensions, shown in
Table~\ref{tab:evaluation-dimensions}.

\begin{longtable}{|p{3.0cm}|p{5.4cm}|p{5.3cm}|}
    \caption{Evaluation dimensions for the proposed system.}
    \label{tab:evaluation-dimensions} \\
    \hline
    \textbf{Dimension} & \textbf{Main question} & \textbf{Typical evidence} \\
    \hline
    \endfirsthead
    \hline
    \textbf{Dimension} & \textbf{Main question} & \textbf{Typical evidence} \\
    \hline
    \endhead

    Functional correctness
        & Do the main instructor, learner, administrator, and LMS workflows
          behave as specified?
        & Test cases, API responses, database state, screenshots, and event
          logs. \\
    \hline

    Integration correctness
        & Do LTI, Deep Linking, AGS, queues, and worker interactions preserve
          the intended contracts?
        & Launch traces, JWT claims, BullMQ job logs, outbox records, and
          Canvas sandbox observations. \\
    \hline

    Retrieval quality
        & Does the system retrieve the right chunks, cards, and transcript
          segments under course, role, and publish-status filters?
        & Ground-truth query set, Recall@K, MRR, nDCG@K, and latency logs. \\
    \hline

    Generated-content quality
        & Are AI-generated learning artifacts correct, clear, useful, and
          aligned with Learning Outcomes?
        & Instructor rubric scores, source-grounding checks, schema validation,
          and rejection reasons. \\
    \hline

    Security and privacy
        & Are authentication, authorization, tenant isolation, prompt-injection
          controls, and audit logs adequate?
        & Security checklist, access-control tests, negative tests, and audit
          samples. \\
    \hline

    Performance and scalability
        & Are user-facing APIs and background pipelines fast enough for a
          demo/pilot class?
        & P50/P95 latency, queue time, processing time, throughput, and resource
          usage. \\
    \hline

    Operations and cost control
        & Can administrators monitor failures, retries, LLM usage, and quota
          limits?
        & Grafana dashboards, LLM usage logs, quota events, alert records, and
          incident runbooks. \\
    \hline
\end{longtable}

\section{Functional Evaluation}

Functional evaluation verifies whether the implemented system satisfies the
main workflows. The tests should be executed as black-box scenarios through
the user interface or public API, while also checking the database and event
state after each scenario.

\begin{longtable}{|p{1.8cm}|p{4.8cm}|p{4.0cm}|p{3.1cm}|}
    \caption{Functional evaluation scenarios.}
    \label{tab:functional-evaluation-scenarios} \\
    \hline
    \textbf{ID} & \textbf{Scenario} & \textbf{Evaluation method} & \textbf{Acceptance criterion} \\
    \hline
    \endfirsthead
    \hline
    \textbf{ID} & \textbf{Scenario} & \textbf{Evaluation method} & \textbf{Acceptance criterion} \\
    \hline
    \endhead

    FE-01
        & Upload a PDF/DOCX/PPTX learning resource.
        & Use an instructor account, upload a valid file, and inspect the
          document record, MinIO object, and queued job.
        & The resource is stored, linked to the course, and enters the
          processing pipeline. \\
    \hline

    FE-02
        & Reject invalid or unsafe uploads.
        & Upload unsupported file types, oversized files, and duplicate files.
        & The system rejects invalid input with a user-friendly error and does
          not create inconsistent storage objects. \\
    \hline

    FE-03
        & Process a document into chunks, vectors, and graph nodes.
        & Run the document pipeline and inspect PostgreSQL, Qdrant, Neo4j, and
          outbox events.
        & Chunks are persisted, vectors are indexed, and graph projection is
          eventually synchronized. \\
    \hline

    FE-04
        & Process a lecture video into transcript segments and indexed video
          content.
        & Upload a short video and inspect video segments, transcript segments,
          generated clips, and vector payloads.
        & The video reaches the \texttt{INDEXED} state and can be retrieved
          through course-scoped search. \\
    \hline

    FE-05
        & Generate lesson cards, flashcards, and quiz items.
        & Trigger content generation for a selected lesson or LO, then inspect
          generated artifacts and schema validation results.
        & Generated artifacts match the required JSON schema and enter the
          review workflow as drafts. \\
    \hline

    FE-06
        & Review, edit, approve, publish, unpublish, and archive generated
          content.
        & Use the instructor review UI and inspect audit logs.
        & Learners only see published content; all state transitions are logged
          and reversible where required. \\
    \hline

    FE-07
        & Learn through lesson cards and flashcards.
        & Use a learner account, open a published lesson, flip cards, and mark
          review state.
        & Card progress and flashcard review state are persisted per learner. \\
    \hline

    FE-08
        & Take a quiz and receive immediate feedback.
        & Submit quiz answers and inspect quiz attempts, score, explanation,
          and optional AGS status.
        & The attempt is stored, feedback is shown, and AGS passback is
          scheduled or marked as not required. \\
    \hline

    FE-09
        & Ask the AI Tutor a course-scoped question.
        & Send questions with and without relevant source material.
        & Answers cite retrieved sources; out-of-scope questions are refused or
          answered as unavailable in the course. \\
    \hline

    FE-10
        & Configure quota and observe usage.
        & Set user, course, and instructor limits, send LLM requests, and inspect
          usage logs.
        & Requests are allowed before quota and rejected or degraded after the
          configured threshold. \\
    \hline
\end{longtable}

\section{LMS and Integration Evaluation}

Canvas LMS integration is evaluated separately because it is a critical
system boundary. The evaluation should use a Canvas sandbox or a controlled
self-hosted Canvas instance.

\subsection{LTI 1.3 Launch}

The LTI launch test verifies the OIDC login flow, JWT verification, role
mapping, and session creation. The evaluator should record:

\begin{itemize}
    \item issuer, client ID, deployment ID, audience, nonce, and expiration
          validation result;
    \item mapped internal user ID and course ID;
    \item course-scoped role from \texttt{course\_memberships};
    \item generated server-side session and short-lived BFF-to-Core JWT;
    \item final redirect path for Instructor, Learner, and Administrator
          contexts.
\end{itemize}

The launch is accepted if the tool rejects invalid signatures, expired tokens,
wrong audiences, replayed nonces, and users without required course roles.

\subsection{Deep Linking}

Deep Linking is evaluated by asking an instructor to select a lesson or
activity from the tool and return it to Canvas. The expected evidence is:

\begin{itemize}
    \item the tool receives a valid Deep Linking request from Canvas;
    \item the selected item is stored in \texttt{lti\_resource\_links};
    \item Canvas creates an activity containing the expected custom claim, for
          example \texttt{custom.lesson\_id};
    \item a learner launch from that activity opens the correct lesson.
\end{itemize}

\subsection{Assignment and Grade Services}

AGS evaluation verifies that a quiz attempt can be mapped to the correct LTI
line item. The test should inspect \texttt{quiz\_attempts},
\texttt{lti\_resource\_links}, and \texttt{lti\_line\_items}. The passback is
accepted if the posted score uses the correct learner ID, score scale, line
item URL, activity progress, and grading progress.

\section{Retrieval Quality Evaluation}

Retrieval quality determines whether RAG and GraphRAG receive the correct
context. It should be evaluated before judging chatbot or quiz-generation
quality, because weak retrieval can make downstream generation fail even when
the prompt is well designed.

\subsection{Ground-Truth Dataset}

The retrieval dataset should contain:

\begin{itemize}
    \item course documents with known chapter and LO structure;
    \item lecture videos with timestamped transcript segments;
    \item 50--100 search queries written by instructors or teaching
          assistants;
    \item ground-truth relevant chunks, lesson cards, or transcript segments
          for each query;
    \item negative queries whose answers should not be available in the course.
\end{itemize}

Each query should be tagged by query type: keyword search, natural-language
paraphrase, LO-based query, video timestamp query, or out-of-scope query.

\subsection{Retrieval Strategies}

The planned evaluation compares four strategies:

\begin{enumerate}
    \item \textbf{Vector-only retrieval:} dense or hybrid Qdrant search without
          graph expansion.
    \item \textbf{Vector + metadata filter:} retrieval with hard filters on
          \texttt{course\_id}, source type, publish visibility, and allowed
          LOs.
    \item \textbf{Graph-aware reranking:} vector retrieval combined with
          Neo4j LO/chapter relationships.
    \item \textbf{GraphRAG for generation:} retrieval constrained by the target
          LO and expanded through prerequisite or related concepts.
\end{enumerate}

\subsection{Retrieval Metrics}

\begin{table}[H]
    \centering
    \caption{Retrieval evaluation metrics.}
    \label{tab:retrieval-evaluation-metrics}
    \begin{tabular}{|p{3.0cm}|p{10.5cm}|}
        \hline
        \textbf{Metric} & \textbf{Meaning} \\
        \hline
        Recall@K
            & Percentage of queries for which at least one relevant item
              appears in the top-K results. \\
        \hline
        Precision@K
            & Percentage of top-K retrieved items that are relevant. \\
        \hline
        MRR
            & Mean Reciprocal Rank of the first relevant item. \\
        \hline
        nDCG@K
            & Ranking quality when relevance is graded as high, medium, low,
              or irrelevant. \\
        \hline
        Refusal accuracy
            & Percentage of out-of-scope queries where the system correctly
              refuses to answer from unsupported context. \\
        \hline
        Retrieval latency P95
            & 95th percentile retrieval time, including embedding and vector
              search. \\
        \hline
    \end{tabular}
\end{table}

\section{Generated Content Quality Evaluation}

Generated content should not be evaluated only by automatic schema validation.
Because the system is used in an educational context, instructor or teaching
assistant review is required.

\subsection{Lesson Card and Flashcard Rubric}

Lesson cards and flashcards are evaluated with the rubric in
Table~\ref{tab:card-rubric}. Each criterion is scored from 1 to 5.

\begin{longtable}{|p{3.7cm}|p{6.7cm}|p{3.1cm}|}
    \caption{Rubric for lesson card and flashcard quality.}
    \label{tab:card-rubric} \\
    \hline
    \textbf{Criterion} & \textbf{Question for reviewer} & \textbf{Pass condition} \\
    \hline
    \endfirsthead
    \hline
    \textbf{Criterion} & \textbf{Question for reviewer} & \textbf{Pass condition} \\
    \hline
    \endhead

    Factual correctness
        & Is the card consistent with the source document or transcript?
        & No serious factual error. \\
    \hline

    Source grounding
        & Can the card be traced to one or more source chunks or transcript
          segments?
        & At least one valid source. \\
    \hline

    Conciseness
        & Is the content short enough for micro-learning while still complete?
        & Average score $\geq 3.5$. \\
    \hline

    Learning usefulness
        & Does the card help the learner understand or remember the target
          concept?
        & Average score $\geq 3.5$. \\
    \hline

    LO alignment
        & Does the card support the intended Learning Outcome?
        & Reviewer agrees or marks minor adjustment only. \\
    \hline
\end{longtable}

\subsection{Quiz Item Rubric}

Quiz items are evaluated on correctness, clarity, Bloom level, distractor
quality, and explanation quality. The review should cover all supported
question types: \texttt{MCQ\_SINGLE}, \texttt{MCQ\_MULTI},
\texttt{TRUE\_FALSE}, and \texttt{FILL\_BLANK}.

\begin{longtable}{|p{3.8cm}|p{6.6cm}|p{3.1cm}|}
    \caption{Rubric for quiz item quality.}
    \label{tab:quiz-rubric} \\
    \hline
    \textbf{Criterion} & \textbf{Question for reviewer} & \textbf{Pass condition} \\
    \hline
    \endfirsthead
    \hline
    \textbf{Criterion} & \textbf{Question for reviewer} & \textbf{Pass condition} \\
    \hline
    \endhead

    Correct answer
        & Is the marked answer correct according to the source?
        & No wrong answer accepted. \\
    \hline

    Question clarity
        & Is the wording clear and not misleading?
        & Average score $\geq 3.5$. \\
    \hline

    Distractor quality
        & For MCQs, are wrong options plausible but clearly incorrect?
        & No obviously absurd distractors. \\
    \hline

    Bloom alignment
        & Does the item match its intended Bloom level?
        & Reviewer agrees or marks one-level adjustment. \\
    \hline

    Explanation quality
        & Does the explanation help learners understand the answer?
        & Average score $\geq 3.5$. \\
    \hline

    Source grounding
        & Is the item supported by source chunks, transcripts, or LOs?
        & At least one valid source. \\
    \hline
\end{longtable}

\subsection{AI Tutor Answer Rubric}

Chatbot answers are evaluated separately from generated cards and quizzes.
For each test question, the reviewer checks:

\begin{itemize}
    \item whether the answer is grounded in retrieved course content;
    \item whether the citations point to the correct document, card, or video
          segment;
    \item whether the answer avoids unsupported claims;
    \item whether the answer is useful for learning rather than merely giving a
          final answer;
    \item whether the system refuses out-of-scope questions correctly.
\end{itemize}

An answer is acceptable if it has no serious factual error, includes at least
one relevant citation when source material exists, and refuses to invent an
answer when no suitable context is found.

\section{Security and Privacy Evaluation}

Security evaluation focuses on the risks introduced by LMS integration,
multi-user course access, generated content, and LLM usage. The planned tests
are listed in Table~\ref{tab:security-tests}.

\begin{longtable}{|p{3.2cm}|p{5.5cm}|p{4.8cm}|}
    \caption{Security and privacy evaluation checklist.}
    \label{tab:security-tests} \\
    \hline
    \textbf{Area} & \textbf{Evaluation method} & \textbf{Acceptance criterion} \\
    \hline
    \endfirsthead
    \hline
    \textbf{Area} & \textbf{Evaluation method} & \textbf{Acceptance criterion} \\
    \hline
    \endhead

    LTI token validation
        & Send expired, replayed, wrong-audience, and wrong-signature launch
          tokens.
        & All invalid launches are rejected. \\
    \hline

    Course isolation
        & Attempt to access another course's documents, lessons, vectors, and
          chatbot context.
        & Cross-course access is denied by both API authorization and retrieval
          filters. \\
    \hline

    Role authorization
        & Execute instructor/admin actions with learner credentials.
        & Unauthorized actions return an error and do not change state. \\
    \hline

    CSRF protection
        & Submit state-changing browser requests without the required CSRF
          token.
        & Requests are rejected except verified LTI launch POSTs. \\
    \hline

    Prompt injection
        & Insert malicious instructions into documents or learner questions.
        & Retrieval and prompts preserve source boundaries; the model is not
          allowed to change system behavior. \\
    \hline

    Audit logging
        & Perform review, publish, delete, quota, and admin actions.
        & Security-relevant actions are recorded with actor, timestamp, entity,
          and change details. \\
    \hline
\end{longtable}

\section{Performance and Reliability Evaluation}

Performance evaluation should be conducted only after the main workflows are
implemented and stable. The purpose is not to optimize prematurely, but to
ensure that the architecture can support a small pilot class and identify
bottlenecks.

\subsection{User-Facing API Metrics}

The following API metrics should be collected under normal pilot load:

\begin{itemize}
    \item P50, P95, and P99 latency for reading lessons, listing cards,
          fetching quizzes, submitting quiz attempts, and opening dashboards;
    \item error rate grouped by endpoint and status code;
    \item response size and pagination behavior for large courses;
    \item chat first-token latency and full-answer latency.
\end{itemize}

\subsection{Pipeline Metrics}

Background pipeline evaluation should record:

\begin{itemize}
    \item document parsing time, chunking time, embedding time, and graph-sync
          time;
    \item video transcription time, segmentation time, clip-generation time,
          and transcript-indexing time;
    \item queue waiting time, retry count, failed-job count, and dead-letter
          count;
    \item LLM provider latency, rate-limit errors, token usage, and cost.
\end{itemize}

\subsection{Reliability Tests}

Reliability tests intentionally inject failures:

\begin{itemize}
    \item stop a worker during document processing and verify retry behavior;
    \item fail a Qdrant or Neo4j operation and verify idempotent recovery;
    \item replay an outbox event and verify that duplicate graph nodes or
          vectors are not created;
    \item simulate LLM provider timeout or HTTP 429 and verify fallback,
          backoff, and user-facing error behavior;
    \item run document deletion twice and verify that the delete saga remains
          idempotent.
\end{itemize}

\section{Planned Test Dataset}

The evaluation dataset should be small enough for manual review but diverse
enough to cover the main system behavior. A recommended pilot dataset is:

\begin{itemize}
    \item 3--5 course documents in PDF/DOCX/PPTX format;
    \item 1 syllabus containing chapters, Learning Outcomes, and assessments;
    \item 2--3 lecture videos with clear audio and timestamped segments;
    \item 50 retrieval queries, including at least 10 out-of-scope queries;
    \item 100 generated quiz items across all supported question types;
    \item 50 generated lesson cards or flashcards;
    \item 20 chatbot questions with expected source references;
    \item at least three Canvas roles: Instructor, Learner, and Administrator.
\end{itemize}

For each dataset version, the team should store the source files, expected
ground-truth annotations, prompt version, retrieval configuration, and model
version. This makes later comparisons meaningful when prompts, models, or
chunking strategies change.

\section{Evaluation Reporting Template}

When the implementation is ready, the final evaluation report should include:

\begin{enumerate}
    \item environment: hardware, Docker Compose configuration, database
          versions, model versions, and test date;
    \item dataset: number of documents, videos, chunks, LOs, generated cards,
          generated quiz items, and test queries;
    \item functional test summary: passed, failed, blocked, and deferred
          scenarios;
    \item retrieval metrics: Recall@K, Precision@K, MRR, nDCG@K, refusal
          accuracy, and latency;
    \item generated-content rubric: average score, pass rate, rejection
          reasons, and representative examples;
    \item performance metrics: API latency, queue time, processing time,
          worker throughput, and resource usage;
    \item security results: passed controls, remaining risks, and mitigation
          plan;
    \item limitations: incomplete features, dataset bias, reviewer bias, and
          external-service constraints.
\end{enumerate}

\section{Chapter Summary}

This chapter defined the evaluation framework for the proposed system. Since
the complete implementation is not yet available, the chapter intentionally
does not claim final measured results. Instead, it specifies the criteria,
datasets, rubrics, metrics, and test procedures needed to evaluate the system
once the implementation reaches pilot readiness. This approach keeps the project
honest while still making the expected evaluation process concrete and
actionable.
